\documentclass[../main.tex]{subfiles}
\graphicspath{{\subfix{../image/}}}

\begin{document}

\chapter{Riemann 度量}\label{chapter:ch13-riemannian-metrics}

在本章中，我们首次将几何学引入光滑流形理论。正如本学科的许多内容一样，我们处理几何的方式是以有限维向量空间理论为模型的。为了在向量空间上定义长度与角度等几何概念，人们使用内积。对于流形而言，恰当的结构是 Riemann 度量，它本质上就是在每一个切空间上选取一个内积，并使之逐点光滑变化。一旦选定了 Riemann 度量，我们便能在光滑流形上定义长度、角度与距离等几何概念。

Riemann 几何本身是一门深刻的学科。要发展出一套完整处理它所需的全部工具，还需要整整一本书的篇幅。（如果你想比我们在这里讨论的更深入地钻研它，可以从\cite{LeeRM}开始，该书对这门学科给出了简明的引论，而你已经具备了其中所需的大部分预备知识。）但我们至少可以介绍它的主要定义和一些重要例子。

在定义 Riemann 度量及与之相关的主要构造之后，我们将说明 Riemann 流形的子流形如何继承诱导的 Riemann 度量。然后我们说明 Riemann 度量如何导出一个距离函数，从而使我们可以把连通的 Riemann 流形视为度量空间。在本章末尾，我们将简要介绍 Riemann 度量的一个推广，称为伪 Riemann 度量，它在 Einstein 的广义相对论中起着核心作用。

\subfile{Chapter-13/section-01.tex}
\subfile{Chapter-13/section-02.tex}
\subfile{Chapter-13/section-03.tex}
\subfile{Chapter-13/section-04.tex}
\subfile{Chapter-13/section-05.tex}

\end{document}
